\section{Finite proof data and an explicit certificate}\label{app:finite-data}
The data contain $2\,577$ stored triples: $670$ comparison parameters
and $1\,907$ linear certificates.  The latter contain $142\,017$
nonzero row-weight terms.  There are at most $610$ variables $t_{j,m}$
for any parameter pair in this domain; their values are not enumerated.
The largest common denominator used by the principal arithmetic check
has $68$ decimal digits.  The minimum verified normalized linear
margin is
\[
 \frac{953977}{42042000000}>0,
\]
attained at $(n,a,k)=(34,19,10)$.  This is a certificate upper-bound
margin after normalization, not a lower bound for an unnormalized
coefficient curvature.

The exact data file \texttt{certificates.json} is identified by the
SHA-256 digest
\begin{center}\small\ttfamily
a77f4fc43aec1b8eecab91c4418b6990\allowbreak
f72dee69d802129c78efd67d2e1bcc5a.
\end{center}
The archived proof data and checker source are mathematical supplements
to the finite theorem; the execution logs alone are not substitutes for
them.

\subsection{An explicit affine comparison at the final parameter gap}
\label{f60:subsec-example}

For $(n,a,k)=(60,33,19)$ the data include a certificate with $193$
nonzero row multipliers, including one assumption row.  Its denominator
is $D_c=10^7$ and its equality-weight numerator is $1395$.
Retaining the assumption row algebraically on the left, rather than
declaring it true, gives the unconditional affine comparison
\begin{equation}
 \Delta_{20}\le\frac{96885639}{100000000}\Delta_{19}
 -\frac{4010142596414187852788421881}{28302876966000000000}.
 \label{f60:eq-example-exact}
\end{equation}
Since the last positive fraction is strictly greater than $141686748$,
\begin{equation}
 \Delta_{20}\le\frac{96885639}{100000000}\Delta_{19}-141686748.
 \label{f60:eq-example-rounded}
\end{equation}
The complete weights are in \texttt{example\_60\_33\_19.json}, which
agrees exactly with the corresponding entry of the master data file.
The comparison has been checked by recomputing both normalization
scales and the rational residual bound.  It applies simultaneously to
all forests with $(n,a)=(60,33)$, rather than to a selected graph.

